\vspace{-1mm}
\section{Experiments and Results}
\label{sec:experiments}
\label{sec:results}

\subsection{Experimental Settings}
\label{sec:experiments:settings}

\subsubsection{Datasets}
\label{sec:experiments:datasets}

We conduct experiments on four datasets: \textbf{Liver} and \textbf{Spine}
for synthetic/phantom evaluation, and \textbf{Arm} and \textbf{Breast} for
real-world evaluation. Liver and Spine are adopted from
Ultra-NeRF~\cite{wysockiUltraNeRFNeuralRadiance2024}: Liver contains 6
CT-based simulated sweeps with 200 frames each at $480\times640$, whereas
Spine contains robotic sweeps of a lumbar-spine phantom with 150 frames each
at $168\times1300$. Arm and Breast are derived from
TUS-REC-2024~\cite{liTUSREC2024} and
TDSC-ABUS-2023~\cite{luoTDSCABUSChallenge2025}, respectively. Arm consists of
freehand forearm scans from 85 volunteers at $480\times640$, while Breast
contains automated 3D breast ultrasound volumes with typical sizes of
approximately $850\times600\times300$ voxels. We use all 6 Liver sweeps, 9
Spine sweeps, 33 Arm sweeps, and 30 publicly available Breast volumes. 
All images are adaptively downsampled to approximately $10^5$ pixels per frame
to balance reconstruction accuracy and GPU memory usage.

\subsubsection{Baseline Methods}
\label{sec:experiments:baselines}

We compare our method with seven representative baseline methods from three
categories: NeRF-based Methods, Ultrasound Gaussian Methods, and Adapted
Gaussian Methods. For NeRF-based Methods, we include
Ultra-NeRF~\cite{wysockiUltraNeRFNeuralRadiance2024} and
NeRF-US~\cite{dagliNeRFUSRemovingUltrasound2024}, both of which
combine implicit neural fields with ultrasound-specific rendering. For
Ultrasound Gaussian Methods, we compare with
UltraGauss~\cite{eidUltraGaussUltrafastGaussian2025},
UltraG-Ray~\cite{duelmerUltraGRayPhysicsBasedGaussian2026}, and
UltraGS~\cite{yangUltraGSRealTimePhysicallyDecoupled2025}, which respectively
emphasize plane-based aggregation, physics-based ray casting, and wave-inspired
Gaussian rendering. For Adapted Gaussian Methods, we adapt vanilla
3DGS~\cite{kerbl3DGaussianSplatting2023} and
GaussianPile~\cite{kongGaussianPileUnifiedSparse2026} to ultrasound
reconstruction by modifying their data interfaces and renderers.
For fair comparison, all baselines are adapted to the same reconstruction
setting and evaluated using the same measured probe poses.

\subsubsection{Evaluation Metrics}
\label{sec:experiments:metrics}

We evaluate reconstruction quality using PSNR, L1, SSIM, LPIPS, DISTS, and
GMSD. Higher PSNR and SSIM indicate better performance, whereas lower L1,
LPIPS, DISTS, and GMSD are preferred. We additionally report training time and
rendering frames per second (FPS) for computational efficiency.

\subsubsection{Implementation Details}
\label{sec:experiments:implementation}

We initialize each scene with $10^6$ Gaussians and optimize it for 30k
iterations using Adam. For AFGM, we set $\rho=1$, use the first retained frame
as $i_0$, adopt the dataset-specific nominal acquisition spacing $d_0$, and use
the predefined canonical scan direction. For the local pathway of PLGR, we use
a shared elevational width $\sigma^{\mathrm{ele}}=1$~mm, a support multiplier
of $2.5$, and $\gamma=0.25$. For the global pathway, we use a support
multiplier of $2.0$, 3-point Gauss--Legendre quadrature with a 2-mm integration
half-window, $\lambda_{\mathrm{ray}}=0.55$, and $T_{\min}=0.35$, with
$\beta=0.8$ for local-global modulation. CGAS is applied every 2k iterations
during the active adaptation stage. Residual-guided densification uses
$q_{\mathrm d}=0.98$, $\tau_{\omega}=10^{-6}$, a per-event candidate budget
of approximately $5\%$ of the current Gaussian population, and a $0.6$-mm
scale threshold. Path-aware pruning uses $\tau_{\alpha}=10^{-4}$,
$q_{\mathrm p}=0.10$, and $q_{\mathrm A}=0.80$. Unless otherwise specified,
the same hyperparameters are used across all sequences without per-sequence
tuning. Rendering FPS is measured on an NVIDIA RTX 4090 using pure forward
passes.

\begin{table*}[!t]

\centering

\caption{Quantitative comparison with baseline methods on four ultrasound
datasets. Panels (a) and (b) report synthetic/phantom and real datasets;
$\uparrow/\downarrow$ indicate higher/lower values are preferred, Time and FPS
are averaged within each panel, and bold and underlined values denote the best
and second-best unrounded results.}

\label{tab:comparison:main}

\tiny

\setlength{\tabcolsep}{1.75pt}

\renewcommand{\arraystretch}{0.68}

\setlength{\aboverulesep}{0.42ex}

\setlength{\belowrulesep}{0.45ex}

\setlength{\cmidrulesep}{0.38ex}

\newcommand{\methodcell}[1]{\rule[-0.45ex]{0pt}{2.30ex}#1}

\resizebox{\textwidth}{!}{%

\begin{tabular}{lrrrrrr|rrrrrr|rr}

\toprule

\multicolumn{15}{c}{\textbf{(a) Synthetic/Phantom Datasets}}\\

\midrule

& \multicolumn{6}{c|}{Liver}

& \multicolumn{6}{c|}{Spine}

& \multicolumn{2}{c}{Efficiency}\\

\cmidrule(lr){2-7}\cmidrule(lr){8-13}\cmidrule(lr){14-15}

Method

& PSNR$\uparrow$ & L1$\downarrow$ & SSIM$\uparrow$ & LPIPS$\downarrow$ & DISTS$\downarrow$ & GMSD$\downarrow$

& PSNR$\uparrow$ & L1$\downarrow$ & SSIM$\uparrow$ & LPIPS$\downarrow$ & DISTS$\downarrow$ & GMSD$\downarrow$

& Time (s)$\downarrow$ & FPS$\uparrow$\\

\midrule

\methodcell{Ultra-NeRF}

&21.68&0.064&0.565&0.453&0.283&0.116

&27.13&0.019&0.684&0.510&0.332&0.224

&475.10&315.20\\

\methodcell{NeRF-US}

&22.07&0.042&\underline{0.597}&0.725&0.487&0.207

&27.46&0.019&0.692&0.555&0.359&0.227

&334.10&344.90\\

\methodcell{UltraGauss}

&27.56&0.030&0.567&0.171&0.142&0.113

&28.25&0.021&0.626&0.227&0.198&0.149

&\textbf{170.60}&\textbf{1269.40}\\

\methodcell{UltraG-Ray}

&27.66&0.029&0.564&\textbf{0.154}&\textbf{0.126}&0.109

&\underline{30.63}&0.016&\underline{0.755}&\underline{0.131}&\textbf{0.130}&\underline{0.109}

&\underline{195.10}&\underline{940.40}\\

\methodcell{UltraGS}

&\underline{28.02}&\underline{0.029}&0.594&0.238&0.201&\underline{0.102}

&30.47&\underline{0.016}&0.737&0.187&0.233&0.119

&386.70&417.50\\

\methodcell{3DGS}

&26.73&0.032&0.572&0.279&0.202&0.128

&27.60&0.022&0.713&0.267&0.255&0.163

&413.50&291.40\\

\methodcell{GaussianPile}

&25.23&0.039&0.468&0.331&0.229&0.165

&28.22&0.020&0.717&0.179&0.215&0.129

&409.80&265.00\\

\methodcell{ULGSplat}

&\textbf{29.03}&\textbf{0.025}&\textbf{0.635}&\underline{0.157}&\underline{0.139}&\textbf{0.099}

&\textbf{31.04}&\textbf{0.014}&\textbf{0.797}&\textbf{0.131}&\underline{0.157}&\textbf{0.098}

&513.35&204.94\\

\midrule

\multicolumn{15}{c}{\textbf{(b) Real Datasets}}\\

\midrule

& \multicolumn{6}{c|}{Arm}

& \multicolumn{6}{c|}{Breast}

& \multicolumn{2}{c}{Efficiency}\\

\cmidrule(lr){2-7}\cmidrule(lr){8-13}\cmidrule(lr){14-15}

Method

& PSNR$\uparrow$ & L1$\downarrow$ & SSIM$\uparrow$ & LPIPS$\downarrow$ & DISTS$\downarrow$ & GMSD$\downarrow$

& PSNR$\uparrow$ & L1$\downarrow$ & SSIM$\uparrow$ & LPIPS$\downarrow$ & DISTS$\downarrow$ & GMSD$\downarrow$

& Time (s)$\downarrow$ & FPS$\uparrow$\\

\midrule

\methodcell{Ultra-NeRF}

&26.37&0.025&0.795&0.274&0.249&0.102

&19.45&0.076&0.569&0.569&0.290&0.162

&450.70&314.70\\

\methodcell{NeRF-US}

&29.43&0.017&0.836&0.348&0.289&0.085

&22.00&0.047&0.650&0.563&0.276&0.153

&474.20&327.10\\

\methodcell{UltraGauss}

&31.86&0.015&0.802&0.201&0.191&0.068

&\underline{27.76}&\underline{0.025}&0.782&\underline{0.190}&\underline{0.180}&0.076

&\textbf{171.50}&\textbf{1321.00}\\

\methodcell{UltraG-Ray}

&33.82&0.013&0.832&0.132&0.139&0.058

&27.16&0.031&0.734&0.272&0.211&0.103

&\underline{192.00}&\underline{1190.30}\\

\methodcell{UltraGS}

&31.47&0.020&0.796&0.199&0.183&0.073

&26.29&0.033&0.720&0.351&0.238&0.103

&441.10&403.10\\

\methodcell{3DGS}

&\underline{33.96}&\underline{0.012}&\textbf{0.889}&\underline{0.109}&\underline{0.135}&\underline{0.045}

&27.08&0.027&\underline{0.793}&0.241&0.200&\underline{0.074}

&460.10&307.40\\

\methodcell{GaussianPile}

&32.63&0.014&0.853&0.153&0.142&0.062

&26.85&0.029&0.749&0.326&0.216&0.103

&438.20&252.70\\

\methodcell{ULGSplat}

&\textbf{34.11}&\textbf{0.011}&\underline{0.887}&\textbf{0.100}&\textbf{0.126}&\textbf{0.045}

&\textbf{28.99}&\textbf{0.023}&\textbf{0.809}&\textbf{0.175}&\textbf{0.159}&\textbf{0.065}

&419.37&285.10\\

\bottomrule

\end{tabular}}

\vspace{-4mm}

\end{table*}


\vspace{-3mm}
\subsection{Comparison with Baseline Methods}
\label{sec:results:comparison}

We compare \methodname{} with seven baseline methods from the NeRF-based,
Ultrasound Gaussian, and Adapted Gaussian categories on the Liver, Spine, Arm,
and Breast datasets, in terms of PSNR, L1, SSIM, LPIPS, DISTS, and GMSD,
together with training time and rendering FPS, as reported in
Table~\ref{tab:comparison:main}, with the corresponding qualitative results
presented in Fig.~\ref{fig:comparison:all_datasets}.
Overall, \methodname{} achieves the best PSNR, L1, and GMSD on all four
datasets, while attaining best or second-best SSIM, LPIPS, and DISTS
throughout, showing consistent pixel-wise accuracy and perceptual fidelity
across anatomical targets and acquisition settings.

\subsubsection{Comparison with NeRF-based Methods}

Compared with Ultra-NeRF and NeRF-US, \methodname{} achieves consistently
higher reconstruction accuracy across all four datasets. In particular,
\methodname{} outperforms the strongest NeRF-based baseline on Liver and
Breast, improving PSNR from $22.07\!\rightarrow\!29.03$ dB and
$22.00\!\rightarrow\!28.99$ dB, respectively. On Spine and Arm,
\methodname{} further raises PSNR from 27.46 to 31.04 dB and from 29.43 to
34.11 dB, while improving SSIM from 0.692 to 0.797 and from 0.836 to 0.887.
It also reduces L1 by 39.9\%, 25.5\%, 33.9\%, and 50.9\% on Liver, Spine,
Arm, and Breast, respectively, with consistent gains in the remaining
perceptual and structural metrics. Across both NeRF-based baselines,
\methodname{} improves all six quality metrics on every dataset while
retaining fast rendering.
Similar trends appear in Fig.~\ref{fig:comparison:all_datasets}:
\methodname{} retains the major echo and shadow patterns in Liver, the layered
and curved responses in Spine, and the principal layer organization and echo
patterns in the real Arm and Breast examples, consistent with its quantitative
advantages over NeRF-based methods.

\subsubsection{Comparison with Ultrasound Gaussian Methods}

Among ultrasound-specific Gaussian methods, \methodname{} achieves the most
consistent reconstruction performance across all four datasets. On Liver,
\methodname{} outperforms UltraGS by 1.01 dB in PSNR
($28.02\!\rightarrow\!29.03$ dB) and improves SSIM from 0.594 to 0.635,
while on Spine it surpasses UltraG-Ray with PSNR increasing from 30.63 to
31.04 dB and L1 decreasing by 10.0\%. On real data, \methodname{} exceeds
UltraG-Ray on Arm by 0.30 dB in PSNR, reduces L1 from 0.013 to 0.011, and
raises SSIM from 0.832 to 0.887; on Breast, it further outperforms UltraGauss
by 1.23 dB in PSNR while improving SSIM from 0.782 to 0.809.
Although UltraG-Ray yields lower LPIPS and DISTS on Liver and lower DISTS on
Spine, \methodname{} achieves the best PSNR, L1, SSIM, and GMSD among the
ultrasound-specific Gaussian methods on all four datasets, with best or
second-best LPIPS and DISTS throughout. Its pure-forward rendering reaches
204.94 FPS on the synthetic/phantom panel and 285.10 FPS on the real-data
panel, retaining fast rendering while prioritizing reconstruction fidelity.
\methodname{} also preserves the major Liver echo and shadow patterns,
layered and curved responses in Spine, and coherent layer organization in the
real Arm and Breast examples in Fig.~\ref{fig:comparison:all_datasets},
consistent with its quantitative advantages over specialized Gaussian
baselines.

\subsubsection{Comparison with Adapted Gaussian Methods}

Compared with vanilla 3DGS and GaussianPile, \methodname{} achieves higher
reconstruction quality across the four datasets, as shown in
Table~\ref{tab:comparison:main}. On Liver, \methodname{} outperforms 3DGS by
2.30 dB in PSNR ($26.73\!\rightarrow\!29.03$ dB) and reduces L1 by 21.4\%,
while on Spine it surpasses the stronger GaussianPile result by 2.82 dB and
decreases L1 from 0.020 to 0.014. On real data, \methodname{} improves 3DGS
on Arm from 33.96 to 34.11 dB while reducing L1 by 3.9\%, although its SSIM
is slightly lower (0.887 versus 0.889). On Breast, it raises PSNR by 1.91 dB
($27.08\!\rightarrow\!28.99$ dB). Overall, \methodname{} improves all six
metrics over GaussianPile on every dataset and all metrics except Arm SSIM
over 3DGS.
Qualitatively, Fig.~\ref{fig:comparison:all_datasets} shows that
\methodname{} retains the major Liver echo pattern, layered Spine responses,
and coherent echo organization in the real examples. Together, these results
show that \methodname{} provides more consistent reconstruction quality than
the generic Gaussian adaptations considered here.

\begin{figure*}[!t]
\centering
\includegraphics[width=\textwidth]{figures/pdf/results_comparison_all_datasets.pdf}
\caption{Qualitative comparisons between \methodname{} and seven baseline
methods on four ultrasound datasets. Yellow boxes highlight representative
regions with magnified views. \methodname{} consistently preserves finer
structures and local echo details than the baselines.}
\label{fig:comparison:all_datasets}
\vspace{-4mm}
\end{figure*}


\vspace{-3mm}
\subsection{Ablation Studies}
\label{sec:results:ablation}

In this section, we conduct ablation studies on the Liver and Arm datasets to
assess the joint and internal contributions of the Acquisition-constrained
Frame Geometry Model (AFGM), Physics-aware Local-Global Renderer (PLGR), and
Contribution-guided Gaussian Adaptation Strategy (CGAS). We first analyze
their joint effects and then examine the key design choices within each
component. The same arrow notation and font formatting are used throughout.

\begin{table}[!t]

\centering

\caption{Joint ablation of AFGM, PLGR, and CGAS.}

\label{tab:ablation:joint}

\normalsize
\setlength{\tabcolsep}{2.2pt}
\renewcommand{\arraystretch}{0.88}

\resizebox{\columnwidth}{!}{%

\begin{tabular}{ccc|rrr|rrr}

\toprule

\multicolumn{3}{c|}{Component}
& \multicolumn{3}{c|}{Liver}
& \multicolumn{3}{c}{Arm}\\

\cmidrule(lr){1-3}\cmidrule(lr){4-6}\cmidrule(lr){7-9}

AFGM & PLGR & CGAS
& PSNR$\uparrow$ & SSIM$\uparrow$ & LPIPS$\downarrow$
& PSNR$\uparrow$ & SSIM$\uparrow$ & LPIPS$\downarrow$\\

\midrule

           &            &            & 25.747 & 0.616 & 0.399
                                      & 21.092 & 0.643 & 0.350\\

\checkmark &            &            & 25.608 & 0.617 & 0.565
                                      & 24.136 & 0.760 & 0.322\\

           &            & \checkmark & 25.795 & 0.616 & 0.387
                                      & 21.398 & 0.650 & 0.340\\

\checkmark &            & \checkmark & 25.595 & 0.617 & 0.573
                                      & 24.153 & 0.760 & 0.322\\

           & \checkmark &            & 29.449 & 0.660 & \underline{0.158}
                                      & 32.785 & 0.865 & 0.112\\

\checkmark & \checkmark &            & \underline{29.648} & \underline{0.667} & 0.183
                                      & \underline{33.263} & \underline{0.872} & \underline{0.108}\\

           & \checkmark & \checkmark & 29.452 & 0.660 & \textbf{0.157}
                                      & 32.817 & 0.866 & 0.112\\

\checkmark & \checkmark & \checkmark & \textbf{29.648} & \textbf{0.667} & 0.182
                                      & \textbf{33.271} & \textbf{0.872} & \textbf{0.108}\\

\bottomrule

\end{tabular}}

\vspace{-6mm}

\end{table}

\subsubsection{Joint Component Analysis}

We first design a $2^3$ factorial ablation in
Table~\ref{tab:ablation:joint} to assess AFGM, PLGR, and CGAS, where
each component is compared across four paired configurations with the other
two fixed. 
Comparing Rows 1/5, 2/6, 3/7, and 4/8 clearly reveals the dominant
contribution of PLGR: relative to its Plane/no-global counterpart, PLGR yields
average PSNR gains of 3.863 and 10.339~dB and LPIPS reductions of 0.311 and
0.223 on Liver and Arm, respectively. In contrast, AFGM and CGAS mainly serve
frame-geometry construction and Gaussian population adaptation rather than
direct fidelity enhancement, so their quality effects are smaller. For
instance, CGAS raises Arm PSNR by 0.032~dB in Rows 5/7, while AFGM lowers Arm
LPIPS by 0.004 in Rows 7/8. 
Thus, PLGR mainly drives fidelity, while AFGM and CGAS provide complementary
geometric and structural support.
When all three are enabled, Row 8 attains the best PSNR/SSIM
29.648/0.667 on Liver and 33.271/0.872 on Arm, demonstrating their
complementary roles in 3D US reconstruction.

\begin{table}[!t]

\centering

\caption{Internal ablation of AFGM, where MP denotes full measured poses,
while MG-S and MG-A retain only step magnitude and axial projection,
respectively.}

\label{tab:ablation:afgm}

\scriptsize
\setlength{\tabcolsep}{2.2pt}
\renewcommand{\arraystretch}{0.88}

\resizebox{0.86\columnwidth}{!}{%

\begin{tabular}{lrrr|rrr}

\toprule

& \multicolumn{3}{c|}{Liver}
& \multicolumn{3}{c}{Arm}\\

\cmidrule(lr){2-4}\cmidrule(lr){5-7}

Variant
& PSNR$\uparrow$ & SSIM$\uparrow$ & LPIPS$\downarrow$
& PSNR$\uparrow$ & SSIM$\uparrow$ & LPIPS$\downarrow$\\

\midrule

MP
& 29.452 & 0.660 & \textbf{0.157}
& 32.817 & 0.866 & 0.112\\

MG-S
& 29.641 & 0.667 & \underline{0.182}
& \underline{33.033} & \underline{0.869} & \underline{0.110}\\

MG-A
& \textbf{29.680} & \textbf{0.669} & 0.185
& 32.779 & 0.865 & 0.126\\

AFGM
& \underline{29.648} & \underline{0.667} & 0.182
& \textbf{33.271} & \textbf{0.872} & \textbf{0.108}\\

\bottomrule

\end{tabular}}

\vspace{-4mm}

\end{table}

\subsubsection{Acquisition-constrained Frame Geometry Model}

Next, we study measured-pose completeness using MP, step-magnitude-only MG-S,
axial-projection-only MG-A, and AFGM without measured geometry, as shown
in Table~\ref{tab:ablation:afgm}. Replacing MP with MG-S
(Rows 1/2) slightly improves PSNR from 29.452 to 29.641 on Liver and from
32.817 to 33.033 on Arm, showing that step magnitude alone is sufficient 
for pixel-wise accuracy.
MG-A (Row 3) improves Liver PSNR/SSIM while remaining comparable to MP on Arm,
suggesting that axial progression captures most useful
geometric information. Consistent with both observations, AFGM preserves the essential
progression cues while discarding measured-pose details, achieving the best
Arm PSNR of 33.271~dB and LPIPS of 0.108 and validating its simplified design.

\begin{table}[!t]

\centering

\caption{Internal ablation of PLGR, where Plane, Fixed, and Adaptive denote
elevational-conditioning variants for local echo formation, while Simple and
Full denote ray-integration variants for global transmission.}

\label{tab:ablation:plgr}

\scriptsize
\setlength{\tabcolsep}{2.2pt}
\renewcommand{\arraystretch}{0.88}

\resizebox{0.975\columnwidth}{!}{%

\begin{tabular}{lrrr|rrr}

\toprule

& \multicolumn{3}{c|}{Liver}
& \multicolumn{3}{c}{Arm}\\

\cmidrule(lr){2-4}\cmidrule(lr){5-7}

Variant
& PSNR$\uparrow$ & SSIM$\uparrow$ & LPIPS$\downarrow$
& PSNR$\uparrow$ & SSIM$\uparrow$ & LPIPS$\downarrow$\\

\midrule

Plane
& 25.092 & 0.600 & 0.427
& 21.438 & 0.676 & 0.344\\

Fixed
& \underline{29.161} & \underline{0.643} & 0.169
& 33.674 & 0.884 & 0.126\\

Adaptive
& \textbf{29.339} & \textbf{0.653} & 0.195
& 33.411 & 0.880 & 0.136\\

Fixed + Simple
& 29.057 & 0.636 & \textbf{0.156}
& \underline{34.024} & \underline{0.887} & \underline{0.113}\\

Fixed + Full
& 29.025 & 0.635 & \underline{0.157}
& \textbf{34.125} & \textbf{0.888} & \textbf{0.112}\\

\bottomrule

\end{tabular}}

\vspace{-6mm}

\end{table}

\subsubsection{Physics-aware Local-Global Renderer}

Subsequently, we examine the rationale of PLGR's designs by
varying elevational conditioning in local echo formation and ray integration
in global transmission, while keeping probe-plane marginalization and
ray-aware association fixed, as shown in Table~\ref{tab:ablation:plgr}.
For local echo formation, Plane omits finite-slice conditioning, while Fixed
and Adaptive use shared and depth-adaptive slice widths, respectively.
Fixed (Rows 1/2) improves PSNR by 4.069 and 12.236~dB on Liver and Arm,
respectively, with lower LPIPS, confirming finite-slice modeling.
Adaptive further improves Liver PSNR/SSIM but
yields worse LPIPS and lower Arm performance, exposing a dataset-dependent
trade-off and favoring Fixed as the simpler, more balanced local formulation.
For global transmission, Simple uses approximate path integration, whereas
Full performs causal integration along the pre-pixel acoustic path.
Simple (Rows 2/4) raises Arm PSNR to 34.024~dB and lowers LPIPS on both
datasets, supporting path-dependent transmission.
Moving from Simple to Full (Rows 4/5) gives a further Arm gain
while leaving Liver results nearly unchanged, suggesting a modest advantage
from complete causal integration. 
These results motivate the final Fixed+Full PLGR, which combines
stable finite-slice conditioning with complete causal transmission and
achieves the strongest Arm performance with comparable Liver accuracy.

\begin{table}[!t]

\centering

\caption{Internal ablation of CGAS, where GD and RGD denote gradient-based and
residual-guided densification, AP and LP denote amplitude-only and local-aware
pruning, and Full CGAS further incorporates path-aware pruning. \#G (k)
denotes the Gaussian population in thousands.}

\label{tab:ablation:cgas}

\scriptsize
\setlength{\tabcolsep}{2.2pt}
\renewcommand{\arraystretch}{0.88}

\resizebox{0.905\columnwidth}{!}{%

\begin{tabular}{lrrr|rrr}

\toprule

& \multicolumn{3}{c|}{Liver}
& \multicolumn{3}{c}{Arm}\\

\cmidrule(lr){2-4}\cmidrule(lr){5-7}

Variant
& PSNR$\uparrow$ & LPIPS$\downarrow$ & \#G (k)
& PSNR$\uparrow$ & LPIPS$\downarrow$ & \#G (k)\\

\midrule

GD + AP
& 29.449 & 0.158 & 628
& \underline{32.785} & \underline{0.112} & 680\\

RGD
& 29.434 & \textbf{0.157} & 1051
& 32.709 & 0.112 & 1023\\

RGD + AP
& 29.443 & 0.158 & 643
& 32.773 & 0.112 & 690\\

RGD + LP
& \textbf{29.454} & 0.158 & 645
& 32.778 & \textbf{0.111} & 691\\

Full CGAS
& \underline{29.452} & \underline{0.157} & 689
& \textbf{32.817} & 0.112 & 713\\

\bottomrule

\end{tabular}}

\vspace{-3mm}

\end{table}

\begin{figure}[!t]
\centering
\includegraphics[scale=0.95]{figures/pdf/results_ablation_afgm.pdf}
\vspace{-3mm}
\caption{Sensitivity analysis of the spacing scale $\rho$ in AFGM.}
\label{fig:ablation:afgm}
\vspace{-6mm}
\end{figure}


\subsubsection{Contribution-guided Gaussian Adaptation Strategy}

Finally, we examine CGAS by contrasting gradient-based densification (GD) with
residual-guided densification (RGD), together with amplitude-only pruning
(AP), local-aware pruning (LP), and path-aware pruning in CGAS, as shown
in Table~\ref{tab:ablation:cgas}. Under the same AP rule (Rows 1/3), replacing
GD with RGD maintains comparable quality.
Without pruning, RGD retains 1051k/1023k Gaussians on Liver/Arm, motivating
population control. AP and LP reduce these populations by
about 39\% and 32\%, respectively, while keeping PSNR/LPIPS nearly unchanged.
Full CGAS retains more Gaussians than LP, indicating that
path-aware pruning preserves globally useful capacity while maintaining stable
quality.

\begin{figure}[!t]
\centering
\includegraphics[scale=0.95]{figures/pdf/results_ablation_plgr.pdf}
\vspace{-3mm}
\caption{Sensitivity analysis of the key PLGR hyperparameters: slice width
$\sigma^{\mathrm{ele}}$, composition weight $\beta$, ray attenuation
$\lambda_{\mathrm{ray}}$, quadrature order $N_q$, and integration half-window
$w$.}
\label{fig:ablation:plgr}
\vspace{-3mm}
\end{figure}


\begin{figure}[!t]
\centering
\includegraphics[scale=0.95]{figures/pdf/results_ablation_cgas.pdf}
\vspace{-4mm}
\caption{Sensitivity analysis of the key CGAS hyperparameters: densification
quantile $q_d$, densification budget $B$, pruning quantile $q_p$, path
protection $q_A$, and amplitude threshold $\tau_\alpha$.}
\label{fig:ablation:cgas}
\vspace{-6mm}
\end{figure}


\vspace{-5mm}
\subsection{Hyperparameter Analysis}

\label{sec:results:hyperparameters}

We further analyze the sensitivity of the key hyperparameters in AFGM, PLGR,
and CGAS on representative Liver and Arm sequences, with all results reported
relative to the corresponding default settings to directly reveal parameter
sensitivity and the rationale of the adopted configuration.

\subsubsection{AFGM Spacing Sensitivity}

In \Cref{fig:ablation:afgm}, we use a dual-dataset response curve to examine
how the spacing scale $\rho$ affects reconstruction under synthetic and real
acquisitions, with $\rho=1$ denoting the nominal acquisition spacing. The
Liver curve favors smaller $\rho$, whereas Arm reaches a shallow optimum
around $\rho=1.25$, only 0.069~dB above the default; this opposite preference
around $\rho=1$ indicates that dataset-specific tuning offers no consistent
cross-domain advantage and therefore supports the nominal setting as a simple,
transferable choice without sequence-specific calibration.

\subsubsection{PLGR Rendering Sensitivity}

To examine the robustness of PLGR to its rendering parameters, we vary the
finite-slice width $\sigma^{\mathrm{ele}}$, composition weight $\beta$, ray
attenuation $\lambda_{\mathrm{ray}}$, quadrature order $N_q$, and integration
half-window $w$ independently, as shown in
\Cref{fig:ablation:plgr}. Among them, $\sigma^{\mathrm{ele}}$ exhibits the
strongest domain-dependent response and $\beta$ introduces a milder Liver--Arm
trade-off, whereas $\lambda_{\mathrm{ray}}$ and $w$ remain stable over their
evaluated ranges, indicating substantially lower sensitivity to attenuation
strength and integration extent. The response to $N_q$ is instead
non-monotonic, with $N_q=2$ causing a clear degradation and the adopted
$N_q=3$ recovering strong performance on both datasets, so the final PLGR
configuration represents a balanced operating point across acquisition
domains rather than a dataset-specific peak.

\subsubsection{CGAS Adaptation Sensitivity}

\Cref{fig:ablation:cgas} summarizes CGAS sensitivity using two heatmaps for
$(q_d,B)$ and $(q_p,q_A)$ and a heat strip for $\tau_\alpha$, jointly showing
quality and Gaussian population changes. Across all configurations, PSNR
varies by less than 0.016~dB and forms broad quality plateaus around the
defaults while the Gaussian population changes noticeably, indicating that
these parameters primarily regulate representation capacity rather than
reconstruction fidelity. Since the adopted settings remain within these stable
regions, they provide robust reconstruction quality with controllable Gaussian
complexity, while avoiding unnecessary sensitivity to individual
densification or pruning hyperparameters.
